This dissertation tested whether a US-trained surface-parking segmentation model can be used to measure off-street surface parking in a British city. The model was applied to 100 km² of Leeds exactly as released, with no UK training data in the primary analysis, and evaluated against 2,037 manually labelled car parks drawn to the source model's own target definition. The error was then decomposed rather than merely reported: attributed exhaustively against independent reference layers, characterised by stratified sampling of 142 image chips adjudicated on the imagery the model actually consumed, and tested by ablation of the building- and road-subtraction stage.

The transfer works asymmetrically. Recall of 0.854 is generally high and less spatially variable than precision, which is 0.571; the model predicts 1.50 times the labelled parking area. Accuracy within the city tracks not distance from the centre but how much parking a cell contains --- an apparent location effect that dissolves once parking abundance is controlled for.

Recognition failure is not the dominant explanation: automated attribution leaves 0.0699 km² of whole-lot FN unresolved, 2.1\% of labelled area and under 3\% of all error area, while sampled adjudication separately identifies reference--imagery disagreement. Boundary placement is the largest single component, at 28.8\% of false-positive and 54.1\% of false-negative area, with the remainder split between confusions with particular look-alike surfaces, disagreement over what counts as parking, and pipeline artefacts. Among the sampled unresolved whole-lot population, irregular arrangement was the most frequently identified mechanism and absent markings among the least. Building and road subtraction is a real trade, buying 0.043 of precision for 0.040 of recall, but it also creates a blind spot of its own: the pre-subtraction output detects rooftop parking better than ground-level parking, and building subtraction deletes four fifths of it.

Under that measured reliability, surface parking covers 3.26\% of the study area as labelled, or 3.30\% after the sampled corrections. The highest shares fall in the inner 2 km, reaching 7.11\% in the 1--2 km band, with a clear decline beyond; and it is strongly concentrated: a median cell gives over 1.71\%, while six exceed 10\%.

\textbf{Contributions.} Four, in ascending order of transferability beyond this case. First, an empirical measurement of how a published US-trained parking segmentation model behaves in a British city, previously unavailable. Second, a reusable method for decomposing area-based segmentation error into boundary effects, attributable confusions, definitional disagreement and pipeline artefacts --- a decomposition that turns a single accuracy figure into a statement about which uses survive. Third, evidence that a correction step justified by a sound premise can create a systematic blind spot, which is an argument for evaluating correction stages rather than assuming them. Fourth, a hold-out test of the spatial grain at which a standard area-calibration factor remains reliable: it holds to about ±7\% at half-city scale within Leeds and fails at the scale of a single square kilometre.

\textbf{Future work.} The calibration factor reduces to a ratio of labelled to predicted area, so testing whether it transfers between cities requires labelled totals over a sample of cells rather than a repeat of the full error analysis undertaken here. This makes a multi-city test more tractable, although its cost and between-city validity remain untested. The typology points to a direction for training data: irregularly arranged car parks rather than simply unmarked ones, on the strength of eleven sampled chips. Imagery carrying a near-infrared band would test the one expected failure this study could not examine. A supplementary experiment (Appendix C) narrows what that validation should compare. Generic fine-tuning improved raw-pixel IoU while trading recall for precision; weighting the loss by positional error categories added neither overall nor selective gain, and its operating points were matched or bettered by adjusting the generic model's threshold alone. The next question
is whether, and which kind of, local supervision earns its labelling cost --- a comparison of zero-shot use, area calibration, generic fine-tuning and visually targeted fine-tuning, judged on locating parking, estimating area, and the precision--recall trade-off each implies.

The Leeds results show that a transferred model cannot be trusted to measure parking area without local validation. Within the same city, local calibration supports area estimation at half-city scale; whether that finding transfers between cities remains open.
